\documentclass[12pt,letterpaper]{article}
\usepackage{graphicx,textcomp}
\usepackage{natbib}
\usepackage{setspace}
\usepackage{fullpage}
\usepackage{color}
\usepackage[reqno]{amsmath}
\usepackage{amsthm}
\usepackage{fancyvrb}
\usepackage{amssymb,enumerate}
\usepackage[all]{xy}
\usepackage{endnotes}
\usepackage{lscape}
\newtheorem{com}{Comment}
\usepackage{float}
\usepackage{hyperref}
\newtheorem{lem} {Lemma}
\newtheorem{prop}{Proposition}
\newtheorem{thm}{Theorem}
\newtheorem{defn}{Definition}
\newtheorem{cor}{Corollary}
\newtheorem{obs}{Observation}
\usepackage[compact]{titlesec}
\usepackage{dcolumn}
\usepackage{tikz}
\usetikzlibrary{arrows}
\usepackage{multirow}
\usepackage{xcolor}
\newcolumntype{.}{D{.}{.}{-1}}
\newcolumntype{d}[1]{D{.}{.}{#1}}
\definecolor{light-gray}{gray}{0.65}
\usepackage{url}
\usepackage{listings}
\usepackage{color}

\definecolor{codegreen}{rgb}{0,0.6,0}
\definecolor{codegray}{rgb}{0.5,0.5,0.5}
\definecolor{codepurple}{rgb}{0.58,0,0.82}
\definecolor{backcolour}{rgb}{0.95,0.95,0.92}

\lstdefinestyle{mystyle}{
	backgroundcolor=\color{backcolour},   
	commentstyle=\color{codegreen},
	keywordstyle=\color{magenta},
	numberstyle=\tiny\color{codegray},
	stringstyle=\color{codepurple},
	basicstyle=\footnotesize,
	breakatwhitespace=false,         
	breaklines=true,                 
	captionpos=b,                    
	keepspaces=true,                 
	numbers=left,                    
	numbersep=5pt,                  
	showspaces=false,                
	showstringspaces=false,
	showtabs=false,                  
	tabsize=2
}
\lstset{style=mystyle}
\newcommand{\Sref}[1]{Section~\ref{#1}}
\newtheorem{hyp}{Hypothesis}

\title{Problem Set 3}
\date{Due: March 25, 2026}
\author{Applied Stats II}


\begin{document}
	\maketitle
	\section*{Instructions}
	\begin{itemize}
	\item Please show your work! You may lose points by simply writing in the answer. If the problem requires you to execute commands in \texttt{R}, please include the code you used to get your answers. Please also include the \texttt{.R} file that contains your code. If you are not sure if work needs to be shown for a particular problem, please ask.
\item Your homework should be submitted electronically on GitHub in \texttt{.pdf} form.
\item This problem set is due before 23:59 on Wednesday March 25, 2026. No late assignments will be accepted.
	\end{itemize}

	\vspace{.25cm}
\section*{Question 1}
\vspace{.25cm}
\noindent We are interested in how governments' management of public resources impacts economic prosperity. Our data come from \href{https://www.researchgate.net/profile/Adam_Przeworski/publication/240357392_Classifying_Political_Regimes/links/0deec532194849aefa000000/Classifying-Political-Regimes.pdf}{Alvarez, Cheibub, Limongi, and Przeworski (1996)} and is labelled \texttt{gdpChange.csv} on GitHub. The dataset covers 135 countries observed between 1950 or the year of independence or the first year forwhich data on economic growth are available ("entry year"), and 1990 or the last year for which data on economic growth are available ("exit year"). The unit of analysis is a particular country during a particular year, for a total $>$ 3,500 observations. 

\begin{itemize}
	\item
	Response variable: 
	\begin{itemize}
		\item \texttt{GDPWdiff}: Difference in GDP between year $t$ and $t-1$. Possible categories include: "positive", "negative", or "no change"
	\end{itemize}
	\item
	Explanatory variables: 
	\begin{itemize}
		\item
		\texttt{REG}: 1=Democracy; 0=Non-Democracy
		\item
		\texttt{OIL}: 1=if the average ratio of fuel exports to total exports in 1984-86 exceeded 50\%; 0= otherwise
	\end{itemize}
	
\end{itemize}
\newpage
\noindent Please answer the following questions:

\begin{enumerate}
	\item Construct and interpret an unordered multinomial logit with \texttt{GDPWdiff} as the output and "no change" as the reference category, including the estimated cutoff points and coefficients.
		
	\item Construct and interpret an ordered multinomial logit with \texttt{GDPWdiff} as the outcome variable, including the estimated cutoff points and coefficients.
	
\end{enumerate}

	\subsection*{Answer Q1.1}	
	I start by viewing the data:
	
	\begin{lstlisting}[language=R]
		View(gdp_data)
\end{lstlisting}

As I realise that  gpd\_data\$GDPWdiff is a numeric variable, I transform it into a character variable with three possible categories: "positive", "negative",
or "no change".\\

However, as differences in GDP between years of 0 are very uncommon, I have decided to code the new variable as follows:

	\begin{lstlisting}[language=R]
	gdp_data <- gdp_data %>%
	mutate(GDPWdiff_char = case_when(
	GDPWdiff > 10 ~ "positive",
	GDPWdiff < -10 ~ "negative",
	TRUE           ~ "no change"
	))
\end{lstlisting}

Now that I have all the required variables, I can construct and interpret an unordered multinomial logit with GDPWdiff\_char as the output and ”no change” as the reference category, including the estimated cutoff points and coefficients.\\

I first define "no change" as the reference category and then code the model:

	\begin{lstlisting}[language=R]
gdp_data$GDPWdiff_char <- relevel(as.factor(gdp_data$GDPWdiff_char), ref = "no change")

model_multinom <- multinom(GDPWdiff_char ~ REG + OIL, data = gdp_data)
\end{lstlisting}

% Table created by stargazer v.5.2.3 by Marek Hlavac, Social Policy Institute. E-mail: marek.hlavac at gmail.com
% Date and time: dc., març 25, 2026 - 19:04:56
\begin{table}[!htbp] \centering 
	\caption{} 
	\label{} 
	\begin{tabular}{@{\extracolsep{5pt}}lcc} 
		\\[-1.8ex]\hline 
		\hline \\[-1.8ex] 
		& \multicolumn{2}{c}{\textit{Dependent variable:}} \\ 
		\cline{2-3} 
		\\[-1.8ex] & negative & positive \\ 
		\\[-1.8ex] & (1) & (2)\\ 
		\hline \\[-1.8ex] 
		REG & 0.907$^{***}$ & 1.276$^{***}$ \\ 
		& (0.199) & (0.192) \\ 
		& & \\ 
		OIL & 0.594$^{**}$ & 0.397 \\ 
		& (0.287) & (0.279) \\ 
		& & \\ 
		Constant & 1.380$^{***}$ & 2.158$^{***}$ \\ 
		& (0.093) & (0.088) \\ 
		& & \\ 
		\hline \\[-1.8ex] 
		Akaike Inf. Crit. & 5,685.957 & 5,685.957 \\ 
		\hline 
		\hline \\[-1.8ex] 
		\textit{Note:}  & \multicolumn{2}{r}{$^{*}$p$<$0.1; $^{**}$p$<$0.05; $^{***}$p$<$0.01} \\ 
	\end{tabular} 
\end{table} 

The results in \textit{Table 1} show the estimate log-odds of a country experiencing negative or positive GDP change relative to the baseline of \texttt{no change}. In short, they suggest that democratic regimes are significantly more likely to experience economic movement, albeit dominated by positive growth (log-odds of 1.276 vs log-odds of 0.907). Oil dependency also moves countries away from stagnation but its associated risk is rather negative, unlike the economic outcomes compared to democracy.
	
	\subsection*{Answer Q1.2}
	
First, I put the categories of gdp\_data\$GDPWdiff\_char in order:

	\begin{lstlisting}[language=R]
gdp_data$GDP_ord <- factor(gdp_data$GDPWdiff_car, 
levels = c("negative", "no change", "positive"), 
ordered = TRUE)
\end{lstlisting}

Then, I run the model with the new variable gdp\_data\$GDPWdiff\_ord:

	\begin{lstlisting}[language=R]
model_ord <- polr(GDP_ord ~ REG + OIL, data = gdp_data, Hess = TRUE)
\end{lstlisting}

This ordered model, as shown in \textit{Table 2}, adds new information on the effects of democracy and oil dependency. When the logical order of the differences in GDP is taken into consideration, democracy becomes the only statistically significant driver of economic change, with a positive effect. Oil dependency does not seem to have a statistically significant effect on GDP change.

% Table created by stargazer v.5.2.3 by Marek Hlavac, Social Policy Institute. E-mail: marek.hlavac at gmail.com
% Date and time: dc., març 25, 2026 - 19:12:55
\begin{table}[!htbp] \centering 
	\caption{} 
	\label{} 
	\begin{tabular}{@{\extracolsep{5pt}}lc} 
		\\[-1.8ex]\hline 
		\hline \\[-1.8ex] 
		& \multicolumn{1}{c}{\textit{Dependent variable:}} \\ 
		\cline{2-2} 
		\\[-1.8ex] & GDP\_ord \\ 
		\hline \\[-1.8ex] 
		REG & 0.439$^{***}$ \\ 
		& (0.073) \\ 
		& \\ 
		OIL & $-$0.152 \\ 
		& (0.114) \\ 
		& \\ 
		\hline \\[-1.8ex] 
		Observations & 3,721 \\ 
		\hline 
		\hline \\[-1.8ex] 
		\textit{Note:}  & \multicolumn{1}{r}{$^{*}$p$<$0.1; $^{**}$p$<$0.05; $^{***}$p$<$0.01} \\ 
	\end{tabular} 
\end{table} 

\section*{Question 2} 
\vspace{.25cm}

\noindent Consider the data set \texttt{MexicoMuniData.csv}, which includes municipal-level information from Mexico. The outcome of interest is the number of times the winning PAN presidential candidate in 2006 (\texttt{PAN.visits.06}) visited a district leading up to the 2009 federal elections, which is a count. Our main predictor of interest is whether the district was highly contested, or whether it was not (the PAN or their opponents have electoral security) in the previous federal elections during 2000 (\texttt{competitive.district}), which is binary (1=close/swing district, 0="safe seat"). We also include \texttt{marginality.06} (a measure of poverty) and \texttt{PAN.governor.06} (a dummy for whether the state has a PAN-affiliated governor) as additional control variables. 

\begin{enumerate}
	\item [(a)]
	Run a Poisson regression because the outcome is a count variable. Is there evidence that PAN presidential candidates visit swing districts more? Provide a test statistic and p-value.

	\item [(b)]
	Interpret the \texttt{marginality.06} and \texttt{PAN.governor.06} coefficients.
	
	\item [(c)]
	Provide the estimated mean number of visits from the winning PAN presidential candidate for a hypothetical district that was competitive (\texttt{competitive.district}=1), had an average poverty level (\texttt{marginality.06} = 0), and a PAN governor (\texttt{PAN.governor.06}=1).
	
\end{enumerate}

\subsection*{Answer Q2.1 \& Q2.2}

I run a Poisson model:

	\begin{lstlisting}[language=R]
model_poisson_mex <- glm(PAN.visits.06 ~ competitive.district + marginality.06 + PAN.governor.06, 
family = poisson(link = "log"), 
data = mexico_elections)
\end{lstlisting}

The results in \textit{Table 3} suggest that there is no solid evidence that PAN presidential candidates visited swing districts more frequently (\$z = -0.477, p = 0.6336)\footnote{I do not know how to show this values in the table, but they are available in the .R file}.\\

Instead, the model indicates that the candidate's strategy was primarily driven by the socio-economic status of the district, with a significant preference for wealthier areas, this is, districts with lower marginality (-2.080 and a highly statistically significant p-value). Therefore, as a district becomes more marginal, the number of presidential visits decreases significantly.\\

The model also suggests that districts in states with a PAN governor received fewer visits than districts in states governed by the opposition. However, this association is only margianlly significant (with a p-value of 0.0617). This might indicate something, but the evidence is not strong enough to dismiss the null hypothesis that districts in states with a PAN governor received more visits than districts in states governed by the opposition.

% Table created by stargazer v.5.2.3 by Marek Hlavac, Social Policy Institute. E-mail: marek.hlavac at gmail.com
% Date and time: dc., març 25, 2026 - 19:28:32
\begin{table}[!htbp] \centering 
	\caption{} 
	\label{} 
	\begin{tabular}{@{\extracolsep{5pt}}lc} 
		\\[-1.8ex]\hline 
		\hline \\[-1.8ex] 
		& \multicolumn{1}{c}{\textit{Dependent variable:}} \\ 
		\cline{2-2} 
		\\[-1.8ex] & PAN.visits.06 \\ 
		\hline \\[-1.8ex] 
		competitive.district & $-$0.081 \\ 
		& (0.171) \\ 
		& \\ 
		marginality.06 & $-$2.080$^{***}$ \\ 
		& (0.117) \\ 
		& \\ 
		PAN.governor.06 & $-$0.312$^{*}$ \\ 
		& (0.167) \\ 
		& \\ 
		Constant & $-$3.810$^{***}$ \\ 
		& (0.222) \\ 
		& \\ 
		\hline \\[-1.8ex] 
		Observations & 2,407 \\ 
		Log Likelihood & $-$645.606 \\ 
		Akaike Inf. Crit. & 1,299.213 \\ 
		\hline 
		\hline \\[-1.8ex] 
		\textit{Note:}  & \multicolumn{1}{r}{$^{*}$p$<$0.1; $^{**}$p$<$0.05; $^{***}$p$<$0.01} \\ 
	\end{tabular} 
\end{table}

\subsection*{Answer Q2.3}

Based on the model I just estimated, the general prediction equation would be:\\

Log-mean of PAN presidential candidate visits = -3.81023 -- 0.08135 * competitive.district -- 2.08014 * marginality.06 -- 0.31158 * PAN.governor.06\\

In the case of the winning PAN presidential candidate for a hypothetical district that was competitive (\texttt{competitive.district}=1), had an average poverty level (\texttt{marginality.06} = 0), and a PAN governor (\texttt{PAN.governor.06}=1), the estimated number of visits is...\\

Log-mean of PAN presidential candidate visits = -3.81023 -- 0.08135 * 1 -- 2.08014 * 0 -- 0.31158 * 1 = -4.20316\\

The estimated number of visits is the result of exponentiating the log-count (the result of the previous calculation):\\

e\^-4.20316 = 0.01495\\

Therefore, \textbf{the estimated number of visits the winning PAN presidential candidate made in that hypothetical district is 0.015}.


\end{document}
